\chapter{Metodologi dan Desain Sistem}

Bab ini menjelaskan rancangan sistem, mulai dari alur keseluruhan, kontrak data, arsitektur kedua model, konfigurasi pelatihan, hingga skenario pengujian. Semua tahapan dijalankan secara lokal pada CPU tanpa layanan komputasi awan.

\section{Gambaran Umum Sistem}

Sistem terdiri atas lima tahap berurutan: penyiapan dataset, pelatihan cGAN, pembangkitan citra, pelatihan classifier evaluasi, dan pengukuran kesesuaian kelas. Alur tersebut ditunjukkan pada Gambar \ref{fig:alur}.

\begin{figure}[H]
\centering
\begin{tikzpicture}[node distance=0.72cm,
  block/.style={rectangle, rounded corners, draw=black, fill=orange!20, text width=9.2cm, minimum height=0.72cm, align=center, font=\small},
  data/.style={rectangle, draw=black, fill=blue!15, text width=9.2cm, minimum height=0.72cm, align=center, font=\small},
  arrow/.style={thick,->,>=stealth}]
\node (d1) [data] {Fashion-MNIST: citra $[B,1,28,28]$ dan label $[B]$};
\node (d2) [data, below=of d1] {Preprocessing: normalisasi ke $[-1,1]$};
\node (t1) [block, below=of d2] {Training cGAN: optimasi D, lalu optimasi G (50 epoch)};
\node (t2) [block, below=of t1] {Pembangkitan citra per kelas dari checkpoint terakhir};
\node (t3) [block, below=of t2] {Pelatihan classifier pada data asli (5 epoch)};
\node (t4) [block, below=of t3] {Pengukuran akurasi data asli dan akurasi generator per kelas};
\draw [arrow] (d1) -- (d2);
\draw [arrow] (d2) -- (t1);
\draw [arrow] (t1) -- (t2);
\draw [arrow] (t2) -- (t3);
\draw [arrow] (t3) -- (t4);
\end{tikzpicture}
\caption{Alur keseluruhan sistem dari data hingga evaluasi}
\label{fig:alur}
\end{figure}

\section{Kontrak Data dan Preprocessing}

Dataset dimuat dengan pustaka \texttt{torchvision} di atas kerangka PyTorch \citep{paszke2019pytorch}. Setiap citra dikonversi menjadi tensor float dan dinormalisasi menggunakan rata-rata dan simpangan baku 0,5 sehingga nilainya berada pada $[-1,1]$, selaras dengan fungsi aktivasi \textit{tanh} pada Generator. Kontrak antar-modul dirangkum pada Tabel \ref{tab:kontrak}.

\begin{table}[H]
\centering
\caption{Kontrak antarmuka antar-modul}
\label{tab:kontrak}
\begin{tabular}{|L{3.2cm}|L{2.2cm}|L{2.6cm}|}
\hline
\textbf{Modul} & \textbf{Masukan} & \textbf{Keluaran} \\
\hline
Dataset & Fashion-MNIST & gambar $[B,1,28,28]\!\in\![-1,1]$, label $[B]$, kelas $0..9$ \\
\hline
Generator & noise $[B,100]$, label $[B]$ & gambar $[B,1,28,28]$ dalam $[-1,1]$ \\
\hline
Discriminator & gambar $[B,1,28,28]$, label $[B]$ & logit $[B,1]$ tanpa sigmoid \\
\hline
\textit{Training loop} & batch gambar dan label & loss D dan loss G per epoch, checkpoint, CSV metrik \\
\hline
Evaluasi & checkpoint, data uji asli & akurasi data asli dan akurasi generator per kelas \\
\hline
\end{tabular}
\end{table}

Setiap fungsi masukan melakukan validasi eksplisit terhadap bentuk tensor, jumlah sampel, perangkat, tipe data label, dan rentang label. Validasi ini mencegah kesalahan senyap ketika \textit{conditioning} label diabaikan atau saat batch berukuran satu.

\section{Arsitektur Generator}

Generator menerima noise $z\in\mathbb{R}^{100}$ dan label kelas $y\in\{0,\dots,9\}$ yang diubah menjadi vektor \textit{one-hot} berdimensi 10. Keduanya digabung menjadi vektor 110 dimensi, lalu diproses melalui dua lapisan \textit{fully connected} berukuran 256 dan 512 dengan aktivasi ReLU, sebelum lapisan keluaran 784 unit dengan aktivasi \textit{tanh}. Keluaran kemudian dibentuk ulang menjadi citra $1\times28\times28$. Konfigurasi lapisan ditunjukkan pada Tabel \ref{tab:generator}.

\begin{table}[H]
\centering
\caption{Konfigurasi lapisan Generator}
\label{tab:generator}
\begin{tabular}{|L{0.6cm}|L{1.9cm}|L{1.5cm}|L{3.5cm}|}
\hline
\textbf{No} & \textbf{Lapisan} & \textbf{Ukuran} & \textbf{Aktivasi} \\
\hline
1 & \textit{Concat} (noise + label) & 110 & -- \\
\hline
2 & \textit{Linear} & 256 & ReLU \\
\hline
3 & \textit{Linear} & 512 & ReLU \\
\hline
4 & \textit{Linear} & 784 & \textit{tanh} \\
\hline
5 & \textit{Reshape} & $1\times28\times28$ & -- \\
\hline
\end{tabular}
\end{table}

\section{Arsitektur Discriminator}

Discriminator meratakan citra menjadi vektor 784 dimensi, menggabungkannya dengan label \textit{one-hot} 10 dimensi, lalu memprosesnya melalui dua lapisan \textit{fully connected} berukuran 512 dan 256 dengan \textit{LeakyReLU} berkemiringan 0,2. Lapisan keluaran menghasilkan satu logit per sampel. Rincian disajikan pada Tabel \ref{tab:discriminator}.

\begin{table}[H]
\centering
\caption{Konfigurasi lapisan Discriminator}
\label{tab:discriminator}
\begin{tabular}{|L{0.4cm}|L{1.9cm}|L{1.7cm}|L{2.0cm}|L{2.0cm}|}
\hline
\textbf{No} & \textbf{Lapisan} & \textbf{Masukan} & \textbf{Keluaran} & \textbf{Aktivasi} \\
\hline
1 & \textit{Concat} flattened image + label & 794 & 794 & -- \\
\hline
2 & \textit{Linear} & 794 & 512 & \textit{LeakyReLU} \\
\hline
3 & \textit{Linear} & 512 & 256 & \textit{LeakyReLU} \\
\hline
4 & \textit{Linear} & 256 & 1 & -- \\
\hline
\end{tabular}
\end{table}

\section{Prosedur Pelatihan}

Setiap iterasi memperbarui Discriminator terlebih dahulu menggunakan sampel asli berlabel 1 dan sampel palsu berlabel 0, lalu memperbarui Generator dengan tujuan menipu Discriminator pada sampel palsunya. Parameter Discriminator dibekukan selama pembaruan Generator. Pseudokode prosedur ini ditunjukkan pada Algoritma \ref{algo:train}.

\begin{algorithm}[H]
 \begin{algorithmic}[1]
    \Procedure{TrainCGAN}{$T, \{(x_i,y_i)\}$, $z_\text{dim}$}
    \State Inisialisasi $G$, $D$, optimizer Adam $O_G$, $O_D$, loss $\ell$
    \For{epoch $= 1$ \textbf{to} $T$}
        \For{setiap batch $(x,y)$}
            \State $z \sim \mathcal{N}(0,I)$; $\tilde{x} \gets G(z,y)$
            \State $\mathcal{L}_D \gets \ell(D(x,y), \mathbf{1}) + \ell(D(\tilde{x},y), \mathbf{0})$
            \State Perbarui $D$ dengan $\nabla_{\theta_D}\mathcal{L}_D$
            \State Bekukan parameter $D$
            \State $\mathcal{L}_G \gets \ell(D(G(z,y),y), \mathbf{1})$
            \State Perbarui $G$ dengan $\nabla_{\theta_G}\mathcal{L}_G$
            \State Aktifkan kembali parameter $D$
        \EndFor
        \State Simpan checkpoint epoch dan metrik loss rata-rata
    \EndFor
 \EndProcedure
 \end{algorithmic}
 \caption{Prosedur pelatihan conditional GAN}
 \label{algo:train}
\end{algorithm}

\section{Konfigurasi Eksperimen}

Konfigurasi pelatihan disimpan dalam berkas YAML agar setiap run dapat direproduksi. Nilai yang dipakai dirangkum pada Tabel \ref{tab:konfigurasi}.

\begin{table}[H]
\centering
\caption{Konfigurasi eksperimen}
\label{tab:konfigurasi}
\begin{tabular}{|L{3.4cm}|L{7.6cm}|}
\hline
\textbf{Parameter} & \textbf{Nilai} \\
\hline
Seed & 42 (Python, NumPy, PyTorch, urutan \textit{shuffle} \textit{DataLoader}) \\
\hline
Ukuran batch & 128 \\
\hline
Jumlah epoch & 50 \\
\hline
Optimizer & Adam, learning rate $2\times10^{-4}$, $\beta=(0{,}5;0{,}999)$ \\
\hline
Fungsi loss & \texttt{BCEWithLogitsLoss} \\
\hline
Dimensi noise & 100 \\
\hline
Jumlah kelas & 10 \\
\hline
Perangkat & CPU \\
\hline
Keluaran & checkpoint per epoch, \texttt{metrics.csv}, grid PNG \\
\hline
\end{tabular}
\end{table}

\section{Skenario Pengujian}

Pengujian dibagi menjadi tiga lapis. Pertama, pengujian unit untuk bentuk keluaran, rentang nilai, aliran gradien, sensitivitas \textit{conditioning}, serta validasi masukan model. Kedua, pengujian kontrak data untuk bentuk, tipe, rentang normalisasi, dan reproduksibilitas urutan \textit{shuffle}. Ketiga, pengujian integrasi training dan evaluasi yang memastikan checkpoint dapat dimuat, metrik tercatat, serta laporan evaluasi dihasilkan. Semua pengujian dijalankan dengan perintah \texttt{pytest tests} menggunakan \texttt{PYTHONPATH=src}.

\section{Rancangan Evaluasi}

Classifier evaluasi adalah MLP dengan satu lapisan tersembunyi berukuran 256 dan dilatih dari data asli selama lima epoch, sehingga tidak bergantung pada bobot \textit{pretrained}. Evaluasi menghasilkan dua angka: akurasi pada data uji asli dan akurasi generator dihitung dari 100 gambar per kelas yang dibangkitkan dengan noise ber-seed tetap. Akurasi generator $a_c$ untuk kelas $c$ dihitung dengan Persamaan (\ref{eq:acc}).

\begin{equation}\label{eq:acc}
a_c = \frac{1}{N}\sum_{i=1}^{N} \mathbb{1}\left[\arg\max_k f_k(G(z_i, c)) = c\right]
\end{equation}

dengan $N=100$ dan $f_k$ adalah logit classifier untuk kelas $k$. Nilai rata-rata akurasi generator adalah rata-rata $a_c$ atas sepuluh kelas.
